% Pista ana-24i-q3 · Anàlisi
% Procedència: PAU Catalunya 2024 · Convocatòria de juny (incidències) · Sèrie 5 · Qüestió 3
\documentclass[11pt,a4paper]{article}
\usepackage{pista}
\pistainfo{Catalunya 2024}{Convocatòria de juny (incidències)}{Sèrie 5}

\begin{document}

\section*{\textcolor{blauPista}{Qüestió 3}}

\noindent Considereu les paràboles $y = f_{a}(x)$, amb $f_{a}(x) = ax^{2} + 2x + 5 - a$, on $a$ és un paràmetre real.

\subsection*{\textit{a)} \normalfont Determineu el valor del paràmetre $a$ per al qual la recta tangent a $y = f_{a}(x)$ en el punt d'abscissa $x = 1$ passa pel punt $(2, 13)$. \hfill\textnormal{[1 punt]}}

\pista{Calcula primer $f_{a}(1)$: observa que, en simplificar, el paràmetre $a$ desapareix i obtens una dada interessant i pràctica. A continuació, calcula $f_{a}'(1)$ (aquí $a$ sí que hi surt) i escriu l'equació de la recta tangent amb la forma punt-pendent, $y - f_{a}(1) = f_{a}'(1)(x - 1)$. Imposa, finalment, que la recta passi pel punt $(2, 13)$ i d'aquesta manera trobaràs el valor d'$a$.}

\subsection*{\textit{b)} \normalfont Calculeu els punts de tall de les paràboles $y = f_{1}(x)$ i $y = f_{3}(x)$. \hfill\textnormal{[0,5 punts]}}

\pista{Iguala $f_{1}(x) = f_{3}(x)$ per trobar les abscisses dels punts de tall. Un cop tinguis els valors de $x$, substitueix-los a qualsevol de les equacions de les paràboles (la que et resulti més còmoda), per a obtenir les ordenades corresponents del punts que et demanen. Recorda que la teva resposta final ha de ser escriure els punts.}

\subsection*{\textit{c)} \normalfont Calculeu l'àrea de la regió situada entre les dues paràboles $y = f_{1}(x)$ i $y = f_{3}(x)$. \hfill\textnormal{[1 punt]}}

\pista{L'àrea entre dues corbes és la integral definida de la diferència de les dues funcions (la de ``dalt'' menys la de ``baix''), entre les abscisses dels punts de tall (que tens de l'apartat anterior). Per a decidir quina paràbola està per sobre, avalua-les totes dues en algun punt intermedi dels extrems d'integració: veuràs ràpidament quina té el valor més gran. Una altra manera que hem estudiat a classe és calcular el valor absolut de la integral definida.}

\end{document}
